\section{A model score that averages over ancestors}
\label{sec:fisher}

A good proposal improves the particle approximation, but a separate question
remains: which score should we compute from those particles? Differentiating
every operation of a chosen finite filter is one answer. Fisher's identity
supplies another, aimed directly at the observed-data score.

Assume the initial, transition, and observation laws have differentiable
densities on a common dominating measure, and that differentiation may pass
through the joint integral. For $x_{0:T}$, the joint density is
\begin{equation}
 p_\theta(x_{0:T},y_{1:T})
 =\mu_\theta(x_0)\prod_{t=1}^T
       f_\theta(x_t\mid x_{t-1})g_\theta(y_t\mid x_t).
 \label{eq:joint-model-density}
\end{equation}
The derivative inside this density holds $x_{0:T}$ fixed. Define its complete
data score by
\begin{equation}
 H_T(x_{0:T})=\nabla_\theta\log\mu_\theta(x_0)
 +\sum_{t=1}^T\left\{
    \nabla_\theta\log f_\theta(x_t\mid x_{t-1})
   +\nabla_\theta\log g_\theta(y_t\mid x_t)\right\}.
 \label{eq:complete-data-score}
\end{equation}
Since $\nabla_\theta p_\theta=p_\theta H_T$, integrating this equality and
dividing by $Z=p_\theta(y_{1:T})$ gives
\begin{align}
 s(\theta)
 &=\frac{\int p_\theta(x_{0:T},y_{1:T})H_T(x_{0:T})\,\dd x_{0:T}}{Z}
 \notag\\
 &=\mathbb E_\theta[H_T(X_{0:T})\mid y_{1:T}].
 \label{eq:fisher-model-score}
\end{align}
The posterior expectation is over whole state histories, including the
initial state. A parameter-dependent stationary initial law contributes the
first term even when no observation is assigned to time zero. In a
latent-innovation representation, the density and the derivative must instead
be derived in those coordinates; the state-density expression cannot be used
unchanged for a singular transition.

The needed history average has a forward recursion. Let $H_t$ be
\eqref{eq:complete-data-score} truncated at time $t$, and define
$A_t(x_t)=\mathbb E[H_t(X_{0:t})\mid X_t=x_t,y_{1:t}]$. Conditional on
$X_t$, the observation at $t$ does not further distinguish its ancestors.
The backward conditional distribution is proportional to
$p_\theta(x_{t-1}\mid y_{1:t-1})f_\theta(x_t\mid x_{t-1})$.
Conditioning the additive score on this distribution gives
\begin{align}
 A_t(x)&=\nabla_\theta\log g_\theta(y_t\mid x)
 \notag\\[-2pt]
 &\quad+\int\left[A_{t-1}(z)
           +\nabla_\theta\log f_\theta(x\mid z)\right]
       p_\theta(z\mid x,y_{1:t-1})\,\dd z.
 \label{eq:fisher-conditional-recursion}
\end{align}
This is the step that a single realized ancestor cannot represent exactly.
It averages all histories compatible with the current state instead of
discarding them when a resampling label is selected.

For a weighted particle approximation, define the backward probabilities
\begin{equation}
 b_t^{ji}=
 \frac{W_{t-1}^j f_\theta(x_t^i\mid x_{t-1}^j)}
      {\sum_kW_{t-1}^k f_\theta(x_t^i\mid x_{t-1}^k)}.
 \label{eq:backward-particle-probabilities}
\end{equation}
All terms are evaluated at the current parameter; the denominator must be
positive. The initial particles must approximate $\mu_\theta$: they may be
drawn from it directly, or drawn from a covering proposal and weighted by
the corresponding initial-law ratio. Initialize
$A_0^j=\nabla_\theta\log\mu_\theta(x_0^j)$ and use
\begin{align}
 A_t^i&=\nabla_\theta\log g_\theta(y_t\mid x_t^i)
       +\sum_j b_t^{ji}\left[A_{t-1}^j
          +\nabla_\theta\log f_\theta(x_t^i\mid x_{t-1}^j)\right],\\
 \widehat s_N&=\sum_i W_T^i A_T^i.
 \label{eq:particle-fisher-recursion}
\end{align}
This is the marginal-score construction in
\citet[Section 2.2, Algorithm 2, equations (20) and (22)]{poyiadjis2011particle},
with the initial state indexed explicitly as time zero. Its relation to
modified differentiation rules is discussed by
\citet[Section 4.2, equations (24)--(25)]{scibior2021dpf}. The displayed
recursion follows directly from the conditional expectation above, with the
initial and observation timing made explicit.

No derivative of the proposal or of the sampled trajectory is missing from
\eqref{eq:particle-fisher-recursion}: its integrand is the complete-data
model score. Proposal design changes how well its expectation is estimated.
This differs from $S_0^N$, where state motion through the executed finite
program is part of the claimed derivative. The Fisher estimator should
therefore be checked against a model-score oracle and checked local density
derivatives. Requiring it to match finite differences of the old OT-reset
scalar would ask it to compute a different quantity.

Particle approximation and normalization still introduce error. With an
exact smoothing draw, $H_T$ is unbiased for $s$ under the stated integrability
conditions; with a finite approximate smoothing distribution,
$\widehat s_N$ is not generally unbiased. Averaging ancestors reduces a
particular source of randomness without proving a universal ranking of full
filters. This remains true when the marginal likelihood estimate itself is
unbiased.

The original Poyiadjis paper gives a useful, specific precedent. Its
Section~2.3 compares the methods against an analytic Gaussian-model score
using 500 particles and 100 replications. The discussion of linear growth in
the marginal method's variance combines a stability conjecture with those
numerical results; Theorem~1 instead proves a quadratic lower bound for the
genealogical construction under its stated mixing and boundedness conditions.
These distinctions prevent a local ancestor-averaging argument from becoming
an unsupported universal variance theorem for the proposed filter.

There is a useful computational overlap with
Section~\ref{sec:model-mixtures}. The transition mixture in the numerator
of \eqref{eq:model-marginal-weights} is also the denominator of
\eqref{eq:backward-particle-probabilities}. The same pairwise densities can
support both the value calculation and the score recursion. A direct
implementation costs $O(TN^2)$, plus the cost of the model derivatives.
Streaming reduces memory use, not the number of pairs. PaRIS is an option
for replacing the exact ancestor sum by a controlled number of backward
draws; its acceptance cost, added variance, and support assumptions must be
measured in the actual model before choosing it over the exact reference.

\subsection{Using both past and future information}

Younis and Sudderth's smoother suggests a complementary proposal design.
Their equation (18) draws from a mixture of forward and backward predictive
approximations, allowing either set of observations to nominate a state.
Their equation (19) retains a generative correction, including an auxiliary
reference density; equation (20) replaces its numerator with a learned
discriminative weight \citep[Section 4]{younis2024}. The distinction matters
here. A learned state predictor is not automatically a model posterior, and
multiplying two posterior approximations can count prior information twice
unless the appropriate reference-density factor is included.

For transition parameters, the score contains functions of
$(X_{t-1},X_t)$, so single-time smoothing marginals alone are insufficient.
A Younis-inspired smoother must supply adjacent-state or path information
and the original model correction. Its most direct role is therefore an
evaluable proposal for the expectation in \eqref{eq:fisher-model-score},
followed by model-based weighting, rather than replacement of that expectation
by the paper's supervised state-prediction loss. A learned proposal would
require its own target-specific training protocol and downstream validation.
